% =============================================================
%  SECTION: Conclusion
%  ASSIGNED TO: [Person Name]
% =============================================================

\newpage
\section{Conclusion}

This project set out to build a complete Molecular Dynamics engine from scratch and use it to
study the structural properties of Krypton across the solid--liquid phase transition. Starting
from first principles and assembling six modules, we arrived at a simulation that
is physically correct, numerically stable, and capable of reproducing experimentally measurable
properties of a real noble gas.

Before trusting any physical results, we checked three simple diagnostics. The relative energy fluctuation during the production phase was $6.99 \times 10^{-5}$, which is well within the $10^{-4}$ limit. The total linear momentum remained close to zero ($\sim 10^{-13}$) throughout, showing that momentum zeroing and Newton’s Third Law are upheld. We also observed no systematic drift in energy in any of the simulations.

The physical analysis of Krypton confirmed the MD engine's accuracy. At $T^* = 0.2$
($T \approx 32.5\,\text{K}$), the simulated $g(r)$ exhibited sharp, well-resolved peaks
matching the clear
signature of long-range crystalline order. At $T^* = 1.0$ ($T \approx 162.5\,\text{K}$),
the $g(r)$ collapsed to a single broad first peak decaying to 1.0 within a few angstroms,
matching the signatures of a liquid with short-range order only.

The melting study drove the system through six temperatures from $T = 5\,\text{K}$ to
$T = 162.5\,\text{K}$. As the temperature increases, $g(r)$ changes from a sharp, regular pattern of FCC peaks at low temperature to a single broad hump at high temperature, with the transition occurring between $T = 97.48\,\text{K}$ and $T = 129.97\,\text{K}$. Bracketing
this transition placed the simulated melting point at $T^*_\text{melt} \approx 0.70$,
corresponding to $T_\text{melt} \approx 114\,\text{K}$. This agrees with the established
Lennard-Jones phase diagram value of $T^*_\text{melt} \approx 0.69$ and is within 2\% of the
experimental melting point of Krypton at atmospheric pressure ($115.79\,\text{K}$). The collapse of long-range order between the temperatures on either side of the transition is consistent with the first-order character of the solid–liquid transition.

% \newpage
Overall, these results show that a Lennard–Jones potential, when implemented carefully, can reproduce the structural and thermodynamic behaviour of a real noble gas with high accuracy. This project also gave hands-on experience with all parts of a molecular dynamics code starting from the interaction model and periodic boundaries to numerical integration, performance optimisation, and analysing structural properties.

% Summarise in prose (no bullet points):
%  1. What was built (the 6-module MD engine).
%  2. How it was validated (energy conservation, zero momentum, fluctuation metric).
%  3. What was observed for Krypton (solid at T*=0.4, liquid at T*=1.0).
%  4. What the melting challenge showed.
%  5. Any limitations or future directions.

% [WRITE HERE — 2–3 paragraphs]
